% Generated draft: requires mathematical and source review.
\documentclass[11pt]{article}
\usepackage[margin=25mm]{geometry}
\usepackage{fontspec}
\setmainfont{Latin Modern Roman}
\usepackage{amsmath,amssymb,mathtools,unicode-math}
\setmathfont{Latin Modern Math}
\usepackage{enumitem,needspace,xurl,hyperref}
\setlength{\parindent}{0pt}
\setlength{\parskip}{5pt}
\setlength{\emergencystretch}{3em}
\newcommand{\N}{\mathbb{N}}
\newcommand{\Z}{\mathbb{Z}}
\newcommand{\Q}{\mathbb{Q}}
\newcommand{\R}{\mathbb{R}}
\newcommand{\C}{\mathbb{C}}
\newcommand{\F}{\mathbb{F}}
\newcommand{\A}{\mathbb{A}}
\newcommand{\PP}{\mathbb{P}}
\newcommand{\E}{\mathbb{E}}
\newcommand{\eps}{\varepsilon}
\newcommand{\dd}{\,\mathrm{d}}
\begin{document}
\section*{UnsolvedMath ID: 30005777}
\textbf{Source problem code:} catalogue entry
\subsection{Definitions and mathematical statement}
The supplied record does not specify the ground field, the braiding, or the precise meaning of rank. The following is a conditional formulation using the standard signed-transposition model; these conventions require confirmation against the report.

Let \(k\) denote the intended ground field, whose characteristic is unspecified. For an integer \(n\geq 2\), let \(S_n\) be the symmetric group on \(\{1,\ldots,n\}\), and let \(V_n\) have basis \(x_{ij}\), \(1\leq i<j\leq n\). Write \(x_{ji}=-x_{ij}\) for reversed indices. Define the quadratic Fomin--Kirillov algebra by
\[
\mathcal E_n(k)=T_k(V_n)/J_n,
\]
where \(J_n\) is the two-sided ideal generated by
\[
\begin{aligned}
 x_{ij}^{\,2}&=0 &&(i\ne j),\\
 x_{ij}x_{kl}-x_{kl}x_{ij}&=0 &&(i,j,k,l\text{ pairwise distinct}),\\
 x_{ij}x_{jk}+x_{jk}x_{ki}+x_{ki}x_{ij}&=0 &&(i,j,k\text{ pairwise distinct}).
\end{aligned}
\]
Here \(T_k(V_n)\) is the tensor algebra. Give \(V_n\) the signed-transposition braiding
\[
 c_n(x_{ij}\otimes x_{kl})
 =x_{s_{ij}(k),s_{ij}(l)}\otimes x_{ij},
 \qquad s_{ij}=(i\ j)\in S_n,
\]
with reversed indices interpreted by the convention above. The Nichols algebra \(\mathcal B(V_n,c_n)\) is the quotient of the braided tensor Hopf algebra \(T_k(V_n)\), whose elements of \(V_n\) are primitive, by its largest homogeneous braided Hopf ideal contained in \(\bigoplus_{d\geq2}V_n^{\otimes d}\).

Under these conventions, the question is whether the canonical graded surjection induced by the identity on \(V_n\),
\[
\pi_n:\mathcal E_n(k)\longrightarrow\mathcal B(V_n,c_n),
\]
is an isomorphism for every \(n\geq2\). Equivalently, does the Nichols defining ideal have no additional relations beyond the displayed quadratic relations?

This formulation uses \(n\) for the number of letters of the symmetric group. The source record does not settle whether its word ``rank'' uses this convention or Coxeter rank \(n-1\), and it supplies no field hypothesis. The intended field must therefore be fixed before this becomes an unconditional statement; no quantification over all fields is asserted here. The record mentions known small-rank cases without specifying their range or hypotheses.
\subsection{Short English statement}
The question asks whether, for every intended symmetric-group size, the quadratic Fomin--Kirillov relations already give the full Nichols algebra for the intended braiding. The displayed formulation uses a standard signed-transposition convention, but the ground field and convention must be confirmed. The supplied record labels the general problem open and mentions unspecified known small-rank cases; these are supplied assessments, not independently verified conclusions.
\subsection{Sources}
\begin{itemize}
\item Supplied local UnsolvedMath registry record: permanent ID \(30005777\), title ``Fomin--Kirillov Algebras as Nichols Algebras,'' source problem code \texttt{catalogue entry}, in the Oberwolfach Reports Open Problems collection.
\item Supplied authoritative reference: \emph{Mini-Workshop: Bridging Number Theory and Nichols Algebras via Deformations} (2024), Oberwolfach Reports, pp.\ 235--272, DOI \texttt{10.4171/owr/2024/5}. Supplied source URL and literature-source link, labelled ``original OWR report'': \url{https://doi.org/10.4171/owr/2024/5}.
\item The supplied original excerpt reads ``The general case is an open problem'' and names the contribution ``Geometric problems related to approximations of Nichols algebras,'' Francesco Esposito, joint work with Giovanna Carnovale and Lleonard Rubio y Degrassi. The record does not separately identify a problem proposer.
\item Literature assessment dated exactly \texttt{2026-08-22}: the supplied curation labels the problem open and states that the literature provides special-case results rather than a general solution. This is supplied dated curation, not independently verified literature assessment. Its background explicitly requests independent verification of time-sensitive status and cited claims.
\item The registry marks the statement \texttt{corrected\_verified} and reports that it was reconstructed and checked against the official Oberwolfach report. This is a supplied verification claim; the report has not been independently consulted for this draft.
\end{itemize}
\end{document}
